%% =============================================================================
%% sm_d9_reversal_map.tex — SM-D, Subsection D.6: Reversal Probability Map
%% Body content only (\subsection level); parent structure in sm_d.tex
%%
%% =============================================================================

\subsection{Reversal probability map}
\label{smd:reversal-map}

\Cref{smd:fig-reversal-map} maps the posterior probability of the poverty
reversal for each state:
$\Pr(\tilde{\alpha}_{\mathrm{pov},s} < 0 \text{ and }
\tilde{\beta}_{\mathrm{pov},s} > 0 \mid \text{data})$, computed from the
unweighted M2 model (block-diagonal SVC).  States with high reversal
probability ($\Pr > 0.9$) are concentrated in the South, Southwest, and
large urban states; states with low probability ($\Pr < 0.1$) are
primarily in the upper Midwest and Mountain West.  The remaining states
($0.1 \le \Pr \le 0.9$) tend to be smaller states where the posterior
intervals are wide due to limited sample sizes.

\begin{figure}[htbp]
  \centering
  \includegraphics[width=0.85\textwidth]{F4_reversal_probability_map.pdf}
  \caption{Choropleth map of the posterior reversal probability from the
    unweighted M2 model (block-diagonal SVC):
    $\Pr(\tilde{\alpha}_{\mathrm{pov},s} < 0 \;\cap\;
    \tilde{\beta}_{\mathrm{pov},s} > 0 \mid \text{data})$.  Darker shading
    indicates higher reversal probability.  The population-average
    probability of reversal is 1.000.}\label{smd:fig-reversal-map}
\end{figure}
